\documentclass[../principal.tex]{subfiles}
\graphicspath{{\subfix{../imagenes/}}}

\begin{document}

\chapter{Funciones}

Las funciones nos permiten escribir bloques de código reutilizables, paramétricos, y abstraer los detalles de las implementaciones de nuestras ideas. Las funciones también son conocidas en inglés como functions.

\section{Definición}

Las funciones permiten agrupar bloques de código para poder utilizarlos y volver a utilizarlos, sin necesidad de escribir el mismo código una y otra vez.

\section{Sintaxis en JavaScript}

Las funciones en JavaScript se declaran usando la palabra reservada $function$, y dándoles un nombre de fantasía.

El siguiente punto a considerar es que después del nombre de la función vienen paréntesis (), donde se definen (o no) los parámetros de la función.

Los parámetros son variables locales a la función que permiten ingresar distintos datos para ser usados, leídos e incluso modificados por la función.

Otra palabra clave importante es la palabra return, que permite que la función emita como resultado un valor al exterior, que puede ser usado por el resto del código.

Por ejemplo, aquí definiremos una función sumar() que tiene dos parámetros, a y b, los cuales son sumados internamente, asignados a una variable local resultado, cuyo valor es luego retornado.

\begin{lstlisting}[language=JavaScript]
// declaración de la función sumar
// con dos parámetros llamados a y b
function sumar(a, b) {
    // declaracion de variable local x
    // cuyo resultado es la suma de a y b
    let resultado = a + b;
    // la función retorna el valor almacenado en resultado
    return resultado;
}
\end{lstlisting}

Para usar una función, tiene que haber sido declarada antes. La declaración anterior nos permite luego sumar números de la siguiente forma.

\begin{lstlisting}[language=JavaScript]
let unaSuma = sumar(3 ,4);
\end{lstlisting}


\section{Funciones en p5.js}

El código en JavaScript que nos da el editor de p5.js cuando creamos un nuevo proyecto es el siguiente:

\begin{lstlisting}[language=JavaScript]
function setup() {
  createCanvas(400, 400);
}

function draw() {
  background(220);
}
\end{lstlisting}

En este código distinguimos cuatro funciones distintas, incluyendo setup(), draw(), createCanvas() y background().

Las funciones setup() y draw() están definidas, pero nunca son usadas de forma explícita. Es trabajo de la biblioteca p5.js el usarlas como las conocemos, haciendo que primero ocurra setup() una vez, y luego ocurra draw() en bucle.

Por otro lado, las funciones createCanvas() y background() están siendo usadas, pero no están declaradas acá, lo están en la biblioteca p5.js.


\end{document}


